\section{Introduction}\label{sec:introduction}
Paired moments occur when a two-component linear mode has one rapidly decaying and one slowly accumulating spectral branch. The zero-temperature condition couples their Laplace values: the slow contribution is not independently prescribed data. The fast exponentials alone are incomplete under a summable reciprocal-rate condition, but this does not by itself settle completeness of the coupled rows. Likewise, density of the range of a moment operator is not injectivity of that operator.

We establish an infinite-dimensional smooth nullspace for a specific coupled operator. The mechanism is the time separation between an interior-support fast transform and a tail starting at the terminal point. A biorthogonal family solves fast moment problems for exponentially decaying sequences; the dependent tail is compact into the required weighted sequence space. A Fredholm intersection with the infinite-dimensional fast-only nullspace gives exact solutions, and one further scalar constraint fixes the endpoint. This argument neither assumes that the compact perturbation is small nor treats the tail as arbitrary.

Our inputs include Cannarsa, Martinez and Vancostenoble \cite{Cannarsa2020}, whose biorthogonal estimate concerns fast exponential families, and Almira's exposition of the Müntz theorem \cite{Almira2007}, which concerns the uncoupled exponential family. Guerrero and Imanuvilov \cite{Guerrero2013} study noncontrollability for a heat equation with memory and some nonzero initial states; their target and initial-state question differ from the zero-start paired moments below. A related rest-controllability paper by Ivanov and Pandolfi is cited in \cite{Guerrero2013}; its original proof was not accessible for this comparison. These observations do not establish chronological novelty or exhaust the literature.
